\documentclass[11pt,a4paper]{article}

\usepackage[margin=2.3cm]{geometry}
\usepackage{amsmath}
\usepackage{booktabs}
\usepackage{enumitem}
\usepackage{hyperref}

\setlength{\parindent}{0pt}
\setlength{\parskip}{0.7em}

\begin{document}

\begin{center}
    {\LARGE \textbf{Extracellular Spike Sorting and Curation}}\\[0.5em]
    {Yunjia Li}
\end{center}

\section{Overview}

This document describes the signal-processing and curation workflow used for high-channel-count extracellular neural recordings. The aim is to provide a repeatable procedure from the raw recordings to a final
set of curated single units. Raw recordings, session-specific results, and experimental figures are intentionally not included.

\section{Data Preparation and Initial Checks}

The recordings are sampled at 30~kHz and stored as 16-bit integer data. Before spike sorting, the recording metadata should be checked to confirm the sampling rate, channel count, recording duration, and voltage conversion factor.

Several utility scripts are included in the repository for these checks.
\texttt{raw\_data\_length.py} is used to verify the recording duration,
\texttt{plot\_raw\_data.py} is used to inspect representative raw channels, and \texttt{spike\_length.py} can be used to compare the duration of processed spike output with the raw recording. If strong periodic or sinusoidal noise is observed, \texttt{psd.py} and \texttt{channel\_spectrum.py} can be used to inspect the frequency content of the signal.

\section{Channel Quality Control}

Ground and reference channels are removed before spike sorting. In the current 128-channel configuration, channels 31, 32, 64, 65, 97, 98, 127, and 128 are therefore excluded.

Channels with impedance greater than or equal to 1~M$\Omega$, or reported as open-circuit or outside the measurable range of the recording system, are also excluded.

\section{Probe Mapping}

The probe mapping must be checked before spike sorting because acquisition channel numbers, binary-file indices, and physical electrode positions.

The custom probe mapping should be verified using the mapping utilities in the repository, including \texttt{shank.py} and, where appropriate,
\texttt{test\_map.py}. For downstream analysis, the Kilosort/Phy probe mapping is treated as the canonical channel assignment.

\section{Spike Sorting}

Spike sorting is performed using Kilosort 4. The raw signal is high-pass
filtered at 300~Hz to isolate extracellular spike activity. Spatial whitening is applied using a local neighbourhood of 32 channels.

Rigid drift correction is used by setting \texttt{nblocks} to 1. The universal threshold (\texttt{Th\_universal}) is set to 9 and the learned threshold (\texttt{Th\_learned}) to 8. Templates are extracted using batches of 60,000 samples, corresponding to 2~s at a sampling frequency of 30~kHz.

After sorting, the Kilosort output should be checked for obvious discontinuities, large unexplained blank/periodic regions, or other signs of data corruption before manual curation.

\section{Manual Curation}

Manual curation is performed in Phy. The main purpose of this step is to correct clusters that have been over-split or under-split during automated sorting.

Clusters showing clearly separated populations in principal-component space may be split. Neighbouring clusters with similar waveforms, spatial locations, and complementary cross-correlograms may be considered for merging.

\section{Quantitative Unit Quality Control}

To reduce the subjectivity of manual curation, al candidate single units are evaluated using objective quality metrics calculated with SpikeInterface.

A unit is retained in the final SUA set only if all of the following criteria are satisfied:
\begin{align}
\mathrm{ISI\ violation\ ratio} &< 0.5,\\
\mathrm{amplitude\ cutoff} &< 0.1,\\
\mathrm{presence\ ratio} &> 0.8,\\
|\mathrm{amplitude\_median}| &> 20~\mu\mathrm{V}.
\end{align}

The \texttt{amplitude\_median} value refers to the median extracellular spike amplitude measured from the recording in physical voltage units. It should not be confused with Kilosort/Phy amplitude variables \texttt{amp} and template scaling \texttt{Amplitude}.

For the current Blackrock setup, the voltage conversion factor is
0.25~$\mu$V/bit. This value must be changed if a different acquisition system or gain is used.

\section{Electrophysiological Characterisation}

After final unit selection, the following features can be extracted to
characterise the isolated units:

\begin{itemize}
    \item spike peak-to-trough amplitude;
    \item signal-to-noise ratio;
    \item trough-to-peak duration;
    \item peak-trough ratio;
    \item mean firing rate;
    \item autocorrelogram rise time.
\end{itemize}

The corresponding scripts are provided in the analysis portion of the repository. Each metric should be interpreted separately from the
quality-control thresholds used for SUA selection.

\section{Workflow Summary}

\begin{verbatim}
Raw recording
      |
      v
Recording integrity and signal inspection
      |
      v
Ground/reference and impedance-based channel exclusion
      |
      v
Probe mapping verification
      |
      v
Kilosort 4 spike sorting
      |
      v
Manual curation in Phy
      |
      v
Objective unit-quality filtering
      |
      v
Electrophysiological characterisation
      |
      v
Final curated SUA/MUA output
\end{verbatim}

\section{References}

\begin{enumerate}
    \item Kilosort: \url{https://github.com/MouseLand/Kilosort}
    \item Phy: \url{https://github.com/cortex-lab/phy}
    \item Phy documentation:
    \url{https://phy.readthedocs.io/en/latest/clustering/}
    \item SpikeInterface:
    \url{https://spikeinterface.readthedocs.io/}
    \item Wu et al. (2026):
    \url{https://doi.org/10.1038/s41467-026-71443-7}
\end{enumerate}

\end{document}
